\chapter{Ausblick auf die Masterarbeit}

In der Masterarbeit ist eine Evaluation der Cyber Range als Forschungsumgebung anhand
reproduzierbarer Cyber-Resilienz-Use-Cases vorgesehen. Der Kern der Evaluation ist ein
vergleichbarer Aufbau: Für jeden Use Case wird ein klar definierter Ausgangszustand
festgelegt und über wiederholte Läufe identisch
bereitgestellt. Auf dieser Basis sollen unterschiedliche Werkzeuge oder Werkzeugkonfigurationen
unter denselben Rahmenbedinungen gegenübergestellt werden. Dabei wird pro Vergleichslauf
jeweils nur die Werkzeugkomponente variiert, während Szenarioablauf, Datenquellen und
Messpunkte konstant bleiben.

Um die Vergleichbarkeit sicherzustellen, wird pro Use Case vorab festgelegt, welche
Beobachtungen als erwartete Ergebnisse gelten und wie diese überprüft werden.
Praktisch bedeutet das, dass für jeden Szenariolauf eine definierte Menge an
sicherheitsrelevanten Ereignissen bzw.\ Artefakten erzeugt wird (z.\,B.\ Port-Scan, Login-Versuche,
unerwartete Prozessstarts oder verdächtige Netzwerkverbindungen) und anschließend geprüft wird,
ob und in welcher Form die eingesetzten Werkzeuge diese Aktivitäten erkennen und sichtbar machen.
Die Auswertung kann dabei u.\,a. umfassen, (1) welche Events/Indikatoren vom Werkzeug tatsächlich
gemeldet wurden, (2) wie schnell eine Meldung nach dem Ereignis erfolgt, (3) wie viele irrelevante
Meldungen (False Positives) entstehen und (4) ob die bereitgestellten Kontextinformationen
für eine Interpretation und Nachvollziehbarkeit ausreichen (z.\,B.\ betroffener Host, Benutzer,
Zeitpunkt, technische Details). Durch wiederholte Läufe desselben Use Cases lässt sich zudem prüfen,
ob die Ergebnisse stabil reproduzierbar sind oder ob es zu signifikanten Abweichungen kommt.


Der Forschungsbeitrag entsteht dadurch, dass eine wiederholbare Experimentierumgebung geschaffen wird,
in der definierte Ereignisse reproduzierbar erzeugt werden können. Dadurch wird es möglich, Werkzeuge
nicht anhand einzelner Beobachtungen, sondern anhand wiederholter, vergleichbarer Läufe zu bewerten
(z.\,B.\ welche Aktivitäten erkannt werden, wie konsistent die Ergebnisse sind und welche Unterschiede
zwischen Werkzeugen auftreten).

